\section{Threshold orders and descent on higher models}
\label{orders:section}

We now identify the extension of the specified volume line with the
moduli trace defined by log canonical thresholds.  The calculation is
local on the base, but its compatibility with pullback will give the
assertion on every higher modification.

We use a slightly more general setup than the original fibration.
Let \(p:W\to S\) be a proper connected-fiber map between smooth compact
K\"ahler manifolds, let \(U\subset S\) have SNC complement, and suppose
that over \(U\) the family is smooth with effective relatively SNC
boundary whose coefficients lie in \([0,1]\).  Let \(D_W\) be a
rational divisor extending that boundary, with
\begin{equation}
 K_{W/S}+D_W\sim_{\Q}p^*Q_S
 \label{orders:adjoint}
\end{equation}
and a specified bundle identification in degree \(m\).
The coefficients on vertical components of \(D_W\) are arbitrary
rational numbers; in particular, they need not be at most one.
We use the thresholds of the sub-pair \((W,D_W)\) to define
\[
 B_S=\sum_P(1-t_P)P,\qquad M_S=Q_S-B_S.
\]
All source modifications below carry the crepant transform of
\(D_W\).  The original setup satisfies these conditions.  The
projective comparison in Section~\ref{compare:section} will use the
allowance of arbitrary vertical coefficients.

\subsection{Thresholds on a resolved transverse disk}
\label{orders:thresholds}

\begin{lemma}
\label{orders:threshold}
Fix a prime divisor \(P\subset S\).  On a resolution over a general
point of \(P\), arrange joint SNC support for the crepant divisor
and \(p^*P\).  Write \(\delta_E\) for the crepant coefficient and
\(a_E=\ord_E(p^*P)>0\) for the components dominating \(P\).  Then
\begin{equation}
                 t_P=\min_{E\mapsto P}\frac{1-\delta_E}{a_E}.
 \label{orders:thresholdformula}
\end{equation}
This minimum computes the threshold over all divisorial valuations,
including valuations exceptional over the chosen resolution.
\end{lemma}

\begin{proof}
For a rational SNC divisor, sub-log-canonicity is equivalent to all
coefficients being at most one.  To recall why this tests further
valuations, in coordinates for its support write the logarithmic
top form
\[
 \frac{dx_1}{x_1}\wedge\cdots\wedge\frac{dx_r}{x_r}
 \wedge dx_{r+1}\wedge\cdots\wedge dx_n.
\]
At any prime on a smooth model its pullback has at most a simple
pole: all singular normal terms in the logarithmic differentials
are multiples of the same normal differential \(dz/z\), so at
most one can occur in the wedge.  If \(b_i\) are the coordinate
orders there, this gives the logarithmic Jacobian inequality.
For coefficients \(c_i\leq1\), the resulting log discrepancy is at
least \(\sum_i b_i(1-c_i)\geq0\).  Necessity is already detected
by a component of coefficient greater than one.  Negative
coefficients cause no change in this criterion.

The horizontal coefficients on the resolved model are at most one,
because the sub-pair is log canonical over the good open and their
coefficients are constant.  The coefficient of a vertical component
in \(D_W+t\,p^*P\) is \(\delta_E+ta_E\).  The preceding criterion
therefore gives exactly \eqref{orders:thresholdformula}.
\end{proof}

At a general smooth point of \(P\), choose base coordinates
\((t,z_2,\ldots,z_b)\) with \(P=(t=0)\), and hold the last \(b-1\)
coordinates fixed.  First resolve the joint divisor on the source.
Properness and generic smoothness of the finitely many relevant
strata allow the chosen point to avoid their bad images.  The
resulting transverse slice has a smooth source, and restriction of
the joint divisor is SNC.  Adjunction of the other base coordinates
changes none of the orders in
\eqref{orders:thresholdformula}.  Components mapping to higher
codimension subsets of \(S\) are absent from this general slice.
We may therefore perform the order calculation on a disk.

Throughout that calculation, a relative canonical form means a
section of
\(\omega_W\otimes p^*\omega_{\Delta_t}^{-1}\).
Wedging with the chosen frame \(dt\) identifies it with a total
canonical form.  This convention remains valid at the central
fiber; it does not presume that \(p\) is a submersion there.

\begin{lemma}[Disk models for the extension test]
\label{orders:diskmodels}
Fix the local root family of Section~\ref{hodge:rootconstruction}
and a resolved transverse disk downstairs.  For every sufficiently
divisible power substitution \(t=u^N\), there is a smooth
equivariant K\"ahler model over \(\Delta_u\) with the following
properties:
\begin{enumerate}[label=\textup{(\roman*)}]
\item
It retains every component of the generic root cover and maps to
the downstairs order resolution.  On the punctured disk it is a
modification of the pulled-back root family.
\item
The horizontal removed divisor \(H_N\), defined by the reduced full
inverse image of the original removed divisor on the punctured
disk and then closed up, is SNC.  Its support and the full central
fiber have joint SNC support.
\item
Every stratum of \((Z_N,H_N)\) dominates the disk and is smooth
over a smaller punctured disk.
\item
For each vertical prime \(E\) on the downstairs order resolution,
some prime on \(Z_N\) dominates \(E\).
\end{enumerate}
Consequently \eqref{hodge:FF} applies to this model, and its top
eigenline tests the extension of the original line as in
Lemma~\ref{hodge:birational}.
\end{lemma}

\begin{proof}
Take the main parts of the fiber products of the root cover, the
downstairs order resolution, and \(\Delta_u\).  Here main parts
means every component dominating the disk, so no generic root
sheet is discarded.  Normalize the finite covers and take a common
projective equivariant resolution, including embedded resolution
of the horizontal boundary and central fiber.  This gives the
required maps.  Finite covers and projective resolutions preserve
the local K\"ahler models from Section~\ref{sec:prelim}.  A prime
above a downstairs prime exists already on the normalized finite
cover; retaining its strict transform proves \textup{(iv)}.

Only horizontal components belong to \(H_N\).  Joint SNC with the
central fiber shows that none of their intersection components can
be vertical.  Indeed, if such a stratum were contained in the
central fiber, it would be contained in one of its components.  In
SNC coordinates the intersection of the specified horizontal
coordinate hyperplanes cannot be contained in the additional
central-fiber coordinate hyperplane.  Properness now makes the
image of every stratum the whole disk.  Generic smoothness and
shrinking the disk give smoothness of all strata off its origin.

On the punctured disk the comparison is an equivariant
modification of smooth SNC families, with the reduced full
inverse image as boundary.  Lemma~\ref{hodge:birational} identifies
the top eigenline and its extension.  If the new ambient
cohomology is merely quasi-unipotent, use its upper canonical
extension; the summand coming from the unipotent original line
has the same canonical extension there.
\end{proof}

\subsection{The exact order calculation}
\label{orders:localcalculation}

Let \(\sigma\) be a meromorphic relative \(m\)-plurivolume on the
downstairs disk model whose good-fiber divisor is \(-m\) times the prescribed
effective SNC log Calabi--Yau boundary.  It may have arbitrary
vertical orders.  Let \(\Omega\) denote its multivalued total root,
obtained by multiplying the relative root by \(dt\), and put
\[
 l_E=\frac1m\ord_E(\sigma\otimes dt^{\otimes m}),
 \qquad a_E=\ord_E(t).
\]
Resolve the joint support of the divisor of this tensor and the
central fiber.  The root construction gives a section \(\tau\)
of the pure highest eigenline on the punctured disk.  We compare
it with the canonical extension after unipotentization.

\begin{lemma}[Order of a root volume]
\label{orders:order}
The section \(\tau\) is meromorphic in that extended line.  Its
normalized order is
\begin{equation}
 \ord^{\mathrm{Hdg}}_0(\tau)
       =\min_E\frac{l_E+1-a_E}{a_E},
 \label{orders:generalorder}
\end{equation}
where \(E\) runs through the central-fiber primes on the downstairs
resolution.  The normalization divides an ordinary order after
\(t=u^N\) by \(N\).  At the displayed order, the corrected section
generates the extension after clearing its denominator.
\end{lemma}

\begin{proof}
For a rational number \(b\), take \(N\) sufficiently divisible to
clear all exponents being used and to make the original local
monodromy unipotent.  On the model of
Lemma~\ref{orders:diskmodels}, test \(u^{-Nb}\tau\), with \(\tau\)
now pulled back as a relative form.  The total form being tested is
\begin{equation}
 (u^{-Nb}\tau)\wedge du
       =\frac{u}{N}\,q^*(t^{-1-b}\Omega),
 \label{orders:changevariables}
\end{equation}
where \(q\) maps the resolved composite cover to the downstairs
model.  This identity is simply \(dt=Nu^{N-1}du\), and includes
the base ramification Jacobian.

Set
\[
                  \alpha_E=l_E+1-a_E(1+b).
\]
Suppose first that every \(\alpha_E\) is nonnegative.
Along horizontal primes the divided order is at least \(-1\),
because the good-fiber coefficients lie in \([0,1]\).
Thus \(t^{-1-b}\Omega\) has logarithmic bound on the entire
downstairs SNC chart.  More explicitly, in a logarithmic
differential basis its coefficient is a unit times a product of
nonnegative rational powers of boundary coordinates.  On a cover
making the form single-valued those factors are bounded.  Pullback
of the logarithmic basis has at most a simple pole along each
prime, by the wedge argument in
Lemma~\ref{orders:threshold}.  Hence its pullback has at most
logarithmic poles.

The factor \(u\) in \eqref{orders:changevariables} removes every
vertical logarithmic pole, since \(u\) has positive integral order
on each central-fiber prime.  Along horizontal primes, the tested
relative form has the required logarithmic behavior by
Lemma~\ref{hodge:birational}; the base factor is a unit on the
punctured disk.  It follows that the tested form belongs to
\[
             (g_N)_*\omega_{Z_N/\Delta_u}(H_N).
\]
By \eqref{hodge:FF} and Lemma~\ref{orders:diskmodels}, it extends
in the original canonical highest line.

Conversely, suppose \(\alpha_E<0\) for one downstairs prime \(E\).
Choose a prime above it that dominates it, and let \(e\) be the
ramification index.  Shifted orders of a pulled-back total form
multiply by \(e\).  Also
\(\ord(u)=ea_E/N\) there, since \(t=u^N\).
The shifted order of \eqref{orders:changevariables} is therefore
\begin{equation}
                   e\alpha_E+\frac{ea_E}{N}.
 \label{orders:failure}
\end{equation}
This is \(e(\alpha_E+a_E/N)\), which is negative for every
\(N>a_E/(-\alpha_E)\), regardless of the dependence of \(e>0\)
on \(N\).  The
tested form then has a vertical pole and fails the extension test.
Such a prime survives on a common resolution, so an unfavorable
valuation cannot be lost by choosing a different resolved model.

For completeness, these two tests give an exact order, not just
one inequality.  First choose an integer \(b\) sufficiently negative
that all \(\alpha_E\) are nonnegative, and then choose a fixed
unipotent power \(N_0\) divisible enough for that first test.
That test shows that a power of \(u\) times the pulled-back \(\tau\) is
holomorphic.  Thus \(\tau\) is meromorphic in the canonical line.
Let \(v\) be its ordinary order after that substitution, and put
\(\beta=v/N_0\).  Under a further power \(N=N_0a\) the order is
\(av\), by Lemma~\ref{hodge:purerules}.  For \(Nb\in\Z\), the
extension test is therefore equivalent to \(b\leq\beta\).

The first part proves extension for
\[
 b\leq b_0:=\min_E\frac{l_E+1-a_E}{a_E},
\]
whereas \eqref{orders:failure} proves failure for every rational
\(b>b_0\) after a sufficiently large further power.  Consequently
\(\beta=b_0\), which proves \eqref{orders:generalorder}.  After
clearing its denominator, the section corrected at \(b=b_0\)
has order zero and hence generates the extended line.
\end{proof}

\subsection{The actual moduli extension}
\label{orders:actualextension}

\begin{theorem}[The threshold trace is the volume extension]
\label{orders:trace}
In the setup of \eqref{orders:adjoint}, choose locally a frame
\(s\) of \(\cO_S(mQ_S)\), and let \(\tau\) be its tautological
root as in Section~\ref{hodge:rootconstruction}.  For every prime
\(P\subset S\),
\begin{equation}
                   \ord_P^{\mathrm{Hdg}}(\tau)=t_P-1.
 \label{orders:traceorder}
\end{equation}

More precisely, let \(\pi:\widetilde V\to V\) be a coordinate
power chart making the local pure variation unipotent, and let
\(\overline{\mathcal E}_\chi\) be its extended highest eigenline.
For a sufficiently divisible positive integer \(a\), the open
identification \(s^{\otimes a}\leftrightarrow
\tau^{\otimes am}\) extends to an isomorphism of actual bundles
\begin{equation}
 \pi^*\cO_S(amM_S)|_{\widetilde V}
       \simeq \overline{\mathcal E}_\chi^{\otimes am}.
 \label{orders:tracebundle}
\end{equation}
One integer \(a\) suffices on the compact base.  The resulting
extension rule commutes with pullback to disks generically
meeting \(U\), after further power substitution on the disk.
All assertions allow arbitrary rational coefficients on the
vertical components of \(D_W\).
\end{theorem}

\begin{proof}
For the chosen frame \(s\), the divided order on every crepant
source resolution is \(l_E=-\delta_E\).  Lemmas~\ref{orders:threshold}
and~\ref{orders:order} give
\[
 \ord_P^{\mathrm{Hdg}}(\tau)
   =\min_{E\mapsto P}\frac{1-\delta_E}{a_E}-1=t_P-1.
\]
Over a prime meeting \(U\), the family is smooth and relatively
SNC, so its threshold is one and the same equality is immediate.
The discriminant is therefore supported on the finite SNC bad
divisor and has rational coefficients.

We spell out the extension at intersections of its components.
Choose coordinates \(t_1,\ldots,t_b\) on \(V\) with bad divisor
\(t_1\cdots t_r=0\), and write
\[
                       B_S|_V=\sum_{i=1}^r b_i(t_i=0).
\]
Coefficients \(b_i\) may be negative or zero.  Take the coordinate
power chart \(t_i=w_i^{N_i}\) with \(N_i\) divisible enough to
make monodromy unipotent and \(N_i b_i\) integral.  On its good
open consider
\begin{equation}
                 \gamma=
                 \left(\prod_{i=1}^r w_i^{N_i b_i}\right)\pi^*\tau.
 \label{orders:correctedroot}
\end{equation}
The just-proved disk calculation says that this section has order
zero in \(\overline{\mathcal E}_\chi\) at a general point of every
coordinate divisor.  In fact, on general parallel transverse
disks it extends as a nonvanishing section.

This statement suffices without an a priori multivariable
meromorphicity assertion.  In a local frame of
\(\overline{\mathcal E}_\chi\), let \(g\) be the coefficient of
\(\gamma\) on the good open.  Near a general point of one boundary
component, expand \(g\) as a Laurent series in its normal
coordinate.  The coefficients are holomorphic in the tangential
parameters.  Every negative coefficient vanishes for the dense
set of centers where the disk test applies, hence vanishes
identically.  The same argument applies to \(g^{-1}\), since
the corrected disk section generates.  Thus \(\gamma\) and its
inverse extend across the divisors away from their intersections.
Hartogs extension then extends both across the remaining
codimension-two set.  Their product is still one, so
\(\gamma\) is a frame everywhere.

Choose \(a\) so that \(amM_S\) is an actual integral power and
all \(am b_i\) are integers.  A frame of that power on \(V\) is
\[
                 s^{\otimes a}\prod_{i=1}^r t_i^{am b_i}.
\]
Its pullback corresponds exactly to \(\gamma^{\otimes am}\).
This proves \eqref{orders:tracebundle}.  Notice the sign:
as a meromorphic section of \(mM_S=mQ_S-mB_S\), the original
frame \(s\) has divided order \(-b_i=t_{P_i}-1\).

Changing \(s\) by a unit \(v\) multiplies the local root by a chosen
\(m\)-th root of \(v\) and gives a locally isomorphic root
family.  Choose these roots on a refinement of each overlap by
small coordinate neighborhoods.  They are holomorphic units, and
two choices differ by a locally constant element of \(\mu_m\).
On the \(am\)-th tensor power the change is exactly \(v^a\),
independent of that choice; the possible root-of-unity factors on
triple overlaps also become one.  Thus no root on an entire
overlap is required.  The identifications match the specified
volume identification and its actual bundle cocycle.
Compactness permits finitely many
charts, so one can choose a common positive \(a\).

Finally, a disk map generically meeting \(U\) lifts to each
needed power chart after a further power on the disk: its
boundary coordinates are powers of the parameter times units,
and the units have roots on a small disk.  On that lifted disk,
\eqref{orders:tracebundle} and
Lemma~\ref{hodge:purerules} give the canonical pullback rule.
Nothing in the proof imposed a restriction on vertical
coefficients; only the horizontal good-fiber bounds were used.
\end{proof}

\begin{corollary}[Descent on every higher model]
\label{orders:bdescent}
For the original fibration and the smooth compact K\"ahler model
\(S\) fixed above, every smooth compact K\"ahler modification
\(\nu:S_1\to S\) satisfies
\[
                         M_{S_1}=\nu^*M_S
                         \quad\text{in }\Pic(S_1)\otimes\Q.
\]
No assumption that \(\nu\) be an isomorphism over \(U\) is needed.
\end{corollary}

\begin{proof}
Resolve the main-component family over \(S_1\), using its crepant
divisor \(D_1\), and put
\[
 L_{S_1}=\nu^*L_S,\qquad
 Q_{S_1}=L_{S_1}-K_{S_1}.
\]
The original bundle identification gives
\[
                    K_{W_1/S_1}+D_1
                    \sim_{\Q}p_1^*Q_{S_1}
\]
with its specified degree-\(m\) identification.
Fix a prime \(P\subset S_1\).  Choose a local frame \(s\) of
\(mQ_S\) near a general point of its image.  On the common
isomorphism open it is a section of the same relative volume
line.  Regard it as a meromorphic section of \(mQ_{S_1}\), and
write \(c_P\) for its divided order there.  This interpretation
uses
\begin{equation}
                 Q_{S_1}=\nu^*Q_S-K_{S_1/S}.
 \label{orders:basejacobian}
\end{equation}
In particular, a pulled-back local frame has order
\(c_P=-\ord_P K_{S_1/S}\); the base Jacobian is divided out.

At a general point of \(P\), take a transverse disk whose
punctured part lies in the common good isomorphism open.
After resolving the source jointly with the disk divisor,
the divided order of the relative volume along a component
\(E\) dominating \(P\) is
\begin{equation}
                     l_E=-\delta_E+a_Ec_P,
 \label{orders:higherorders}
\end{equation}
where \(\delta_E\) is the crepant coefficient and
\(a_E=\ord_E(p_1^*P)\).  This is the divisor identity obtained
from the specified relative adjoint isomorphism.  Adjunction of
the other base coordinates introduces no further order.

Use the pulled-back original local root family, and then the
common resolved disk models of
Lemma~\ref{orders:diskmodels}.  Their punctured fibers differ
from that family only by the modifications covered by
Lemma~\ref{hodge:birational}.  Thus
Lemma~\ref{orders:order} computes the order of the original
pulled-back Hodge line and gives
\[
 \min_{E\mapsto P}\frac{l_E+1-a_E}{a_E}
       =c_P+t_{P,S_1}-1.
\]
The right side is precisely the divided order of \(s\) in
\(M_{S_1}=Q_{S_1}-B_{S_1}\).
By Theorem~\ref{orders:trace} and its disk pullback rule, the
same Hodge order is the divided order in \(\nu^*M_S\).
The specified open identification therefore has neither a zero
nor a pole at \(P\).

This argument applies to every prime, including exceptional
primes with image of codimension at least two in the bad
divisor, and primes whose centers lie inside \(U\).
For the latter case one can also check the orders directly:
before further source resolution, smooth base change replaces
the boundary by its pullback minus \(p_1^*K_{S_1/S}\).
If \(j_P=\ord_P K_{S_1/S}\), the threshold is \(1+j_P\);
together with \(c_P=-j_P\) this gives order zero.

To conclude globally in the holomorphic Picard group, the
difference of the two rational lines has the canonical
divisorial description
\[
 M_{S_1}-\nu^*M_S
   =-K_{S_1/S}-B_{S_1}+\nu^*B_S.
\]
The Jacobian supplies the relative canonical divisor even
though the individual canonical bundles need not have global
meromorphic sections.  The order comparison shows that this
rational divisor is zero.  Equivalently, after clearing
denominators the fixed open isomorphism extends across every
divisor and then across codimension two.  This proves the
asserted equality of actual rational line bundles.
\end{proof}

\subsection{The residue integral as an extension norm}
\label{orders:residuenormsection}

We finish with the form of the extension rule needed for the
auxiliary product families.  Let \(r:P\to S\) be a holomorphic
map from a smooth compact complex manifold.  On a dense open
\(P^\circ\) mapping into \(U\), suppose a connected fiber
component \(A_x^*\) of the deepest stratum has been selected
in family.  Write \(s_A\) for the adjunction plurivolume
induced by \(s\in r^*\mathcal G_m\), and define
\begin{equation}
 \|s\|_{\mathrm{res}}
      =\left(\int_{A_x^*}|s_A|^{2/m}\right)^{m/2}.
 \label{orders:residuenormformula}
\end{equation}
For \(k=0\) use the whole connected original fiber, and for a
zero-dimensional selected fiber interpret the integral as
evaluation at its point.

\begin{corollary}
\label{orders:residuenorm}
The expression \eqref{orders:residuenormformula} is a positive
finite norm, homogeneous of degree one, on the volume line.
Let \(a\) clear the indicated powers.  Along any disk in \(P\)
whose punctured part lies in \(P^\circ\), after a sufficiently
divisible power substitution, a nonvanishing local frame of
\[
                        r^*\cO_S(amM_S)
\]
and its dual have at most powers-of-log growth in the \(a\)-th
tensor power of this norm and its dual norm.  No K\"ahler
assumption on \(P\) is required.
\end{corollary}

\begin{proof}
The selected fiber has effective SNC coefficients less than
one.  In local coordinates its measure has factors
\(|z_i|^{-2\delta_i}\), which are integrable for
\(\delta_i<1\).  The integral is therefore finite and positive
for a nonzero volume.  Multiplication by a scalar \(c\)
multiplies it before the outer exponent by \(|c|^{2/m}\),
so the norm has degree one.

Pull the local root family to the punctured disk and retain
all sheets above the selected connected component.  Its
resolved residue cover has a fixed finite degree \(e\).
Proposition~\ref{hodge:residue} gives the specified top-line
identification and
\[
 \int_{\widetilde A_x^*}|\operatorname{Res}\tau|^2
                   =e\int_{A_x^*}|s_A|^{2/m},
\]
up to fixed normalization constants.  Exceptional subsets do
not affect either integral.  On a compact K\"ahler
\((d-k)\)-fold a holomorphic top form is primitive, and its
Hodge norm is this volume integral, again up to a fixed
constant.  Hence \eqref{orders:residuenormformula} is the
\(m\)-th power of the residue Hodge norm, multiplied by
\(e^{-m/2}\).

Theorem~\ref{orders:trace} identifies a pulled-back frame of
the asserted power of \(M_S\) with a canonical highest-line
frame after power substitution.  The pure residue map
preserves such frames by Proposition~\ref{hodge:residue}.
The two-sided bounds in Lemma~\ref{hodge:purerules}, applied
to the compact residue cohomology and then raised to the
required power, give the conclusion.  These are pullbacks
of polarized variations and comparisons of their specified
lines; the smooth compact parameter space itself need not
be K\"ahler.
\end{proof}
